% !TeX root = ../../Book1.tex
\section{光滑逼近}
\label{b1:sec:2201}

\chapref{b1:ch:21}已用局部卷积建立 Sobolev 函数的运算法则。现在要把这种局部逼近接成一个在全域范数下有效的光滑函数。困难不在卷积本身，而在于一个固定半径的核可能越过边界，许多个局部误差又未必能够相加。本节为不同位置选择不同的磨光尺度，同时控制支集和误差。

Meyers 与 Serrin 的短文《\(H=W\)》\cite[pp. 1055--1056]{MeyersSerrin1964HW}证明了任意开集上的光滑稠密性。原文的 \(H^{k,p}\) 指域内光滑函数在 \(W^{k,p}\) 中的闭包，不是本书专用于 \(W^{k,2}\) 的记号 \(H^k\)。以下直接陈述稠密性，不另外引入同名空间。

% 实读来源：MeyersSerrin1964HW，PNAS 51(6) (1964), pp.1055--1056，
% DOI 10.1073/pnas.51.6.1055，PMCID PMC300210；Europe PMC正式扫描副本两页均已阅读。
% p.1055：任意开集、复值、k>=0、1<=p<infty，原文用各导数Lp范数之和；
% p.1056：局部Sobolev函数的全局小误差逼近、环层单位分解及逐层卷积证明。
% 本节用已证完备性处理可和误差级数，补明紧支集补零、支集膨胀与局部有限性；
% 原文末段关于光滑到边界的区别仅作范围说明，不改成本书的紧支集稠密性结论。

\subsection{局部卷积}
\label{b1:subsec:220101}

设 \(\Omega\subseteq\mathbb R^d\) 为非空开集，\(k\in\mathbb N_0\)。仍使用\secref{b1:sec:2003}的非负偶磨光函数 \(\rho\)，其支集包含于 \(\overline B(0,1)\)，积分等于 \(1\)，并记 \(\rho_\varepsilon(x)=\varepsilon^{-d}\rho(x/\varepsilon)\)。对 \(u\in W^{k,p}_{\mathrm{loc}}(\Omega)\)，卷积 \(u*\rho_\varepsilon\) 只先定义在
\[
 \Omega_\varepsilon
 =\{\,x\in\Omega\mid\operatorname{dist}(x,\Omega^c)>\varepsilon\,\}.
\]
由\cref{b1:lem:sobolev-local-smoothing}，在任意固定内部相对紧开集上，这些卷积在有限 \(p\) 的 Sobolev 范数中趋于 \(u\)，并与不超过 \(k\) 阶的弱导数交换。不能在尚未定义 \(u\) 的地方替它指定零值，再假定这些导数关系仍成立。

若函数的非零部分已经留在内部紧集上，补零却没有这个问题。下面将“具有内部紧支集”明确解释为：存在 \(K\Subset\Omega\)，使函数在 \(\Omega\setminus K\) 上几乎处处为零。

\begin{lemma}\label{b1:lem:compact-sobolev-zero-extension}
设 \(1\leq p\leq\infty\)、\(u\in W^{k,p}(\Omega)\)，且 \(u\) 具有内部紧支集 \(K\)。在 \(\Omega\) 外将 \(u\) 补为零，所得函数记为 \(\widetilde u\)。则 \(\widetilde u\in W^{k,p}(\mathbb R^d)\)，并且
\begin{equation}\label{b1:eq:compact-sobolev-zero-extension}
 \partial^\alpha\widetilde u=\widetilde{\partial^\alpha u}
 \quad(|\alpha|\leq k),
 \quad
 \|\widetilde u\|_{W^{k,p}(\mathbb R^d)}
 =\|u\|_{W^{k,p}(\Omega)}.
\end{equation}
\end{lemma}
\begin{proof}
由弱导数的局部性及唯一性，\(\partial^\alpha u=0\) 几乎处处于 \(\Omega\setminus K\)。选取 \(\chi\in C_c^\infty(\Omega)\)，使其在 \(K\) 的一个邻域内等于 \(1\)；这样的截断由\cref{b1:lem:smooth-cutoff}给出。

任取 \(\varphi\in C_c^\infty(\mathbb R^d)\)。函数 \(\chi\varphi\) 是 \(\Omega\) 中的测试函数，而在 \(K\) 附近有 \(\partial^\alpha(\chi\varphi)=\partial^\alpha\varphi\)。因此
\[
 \begin{aligned}
 \int_{\mathbb R^d}\widetilde u\,\partial^\alpha\varphi\,\mathrm dx
 &=\int_\Omega u\partial^\alpha(\chi\varphi)\,\mathrm dx\\
 &=(-1)^{|\alpha|}\int_\Omega\partial^\alpha u\,\chi\varphi\,\mathrm dx\\
 &=(-1)^{|\alpha|}\int_{\mathbb R^d}
        \widetilde{\partial^\alpha u}\,\varphi\,\mathrm dx.
 \end{aligned}
\]
候选导数属于 \(L^p(\mathbb R^d)\)，于是得到弱导数等式。补零不改变任何分量的 \(L^p\) 范数，对有限 \(p\) 求 \(p\) 次和、对无穷端点取最大值，即得范数等式。
\end{proof}

特别地，若只有 \(u\in W^{k,p}_{\mathrm{loc}}(\Omega)\)，而 \(\psi\in C_c^\infty(\Omega)\)，则 \(\psi u\) 仍可使用这个引理。由\cref{b1:prop:weak-leibniz}，
\[
 \partial^\alpha(\psi u)
 =\sum_{\beta\leq\alpha}\binom\alpha\beta
       (\partial^{\alpha-\beta}\psi)(\partial^\beta u).
\]
每个光滑系数都在同一个内部紧集上有界，其支集外为零，所以每项全局属于 \(L^p(\Omega)\)。由此 \(\psi u\in W^{k,p}(\Omega)\)，且具有内部紧支集。这是随后使用单位分解的基本步骤。

\subsection{Friedrichs 磨光函数}
\label{b1:subsec:220102}

Friedrichs 磨光使用上述光滑核把函数换成邻域内的加权平均。\secref{b1:sec:2003}已经解释了这一操作为何使分布光滑；这里还要保留全部 Sobolev 范数估计。

\begin{proposition}\label{b1:prop:sobolev-friedrichs-smoothing}
若 \(u\in W^{k,p}(\mathbb R^d)\)、\(1\leq p\leq\infty\)，则
\(u_\varepsilon=\rho_\varepsilon*u\) 属于 \(C^\infty(\mathbb R^d)\cap W^{k,p}(\mathbb R^d)\)，且
\[
 \partial^\alpha u_\varepsilon
 =\rho_\varepsilon*\partial^\alpha u,
 \quad
 \|u_\varepsilon\|_{W^{k,p}(\mathbb R^d)}
 \leq\|u\|_{W^{k,p}(\mathbb R^d)}.
\]
若 \(p<\infty\)，还有 \(u_\varepsilon\to u\) 于 \(W^{k,p}(\mathbb R^d)\)。
\end{proposition}
\begin{proof}
光滑性和导数交换由\cref{b1:prop:distribution-convolution-smooth}成立。对 \(f\in L^p(\mathbb R^d)\)、\(p<\infty\)，在概率测度 \(\rho_\varepsilon(y)\,\mathrm dy\) 上使用 Hölder 不等式，得到
\[
 \bigl|\rho_\varepsilon*f(x)\bigr|^p
 \leq\int_{\mathbb R^d}\rho_\varepsilon(y)|f(x-y)|^p\,\mathrm dy.
\]
对 \(x\) 积分，Tonelli 定理与平移不变性给出 \(\|\rho_\varepsilon*f\|_p\leq\|f\|_p\)。当 \(p=\infty\) 时，非负核与单位积分直接给出同一估计。将它应用于有限多个导数，便得 Sobolev 范数界。

若 \(p<\infty\)，由\cref{b1:prop:approximate-identity-lp}，每个 \(\rho_\varepsilon*\partial^\alpha u\) 都在 \(L^p\) 中趋于 \(\partial^\alpha u\)。这些误差的 \(p\) 次方只有有限项，其和也趋零，得到所需范数收敛。
\end{proof}

无穷端点的结论只有范数不增，不能把最后一句也扩展到 \(p=\infty\)。\secref{b1:sec:2103}已用 \(|x|\) 的导数跳跃说明：即使在一个有界区间内部，光滑逼近也未必在 \(W^{1,\infty}\) 范数中成立。

\begin{proposition}\label{b1:prop:compact-sobolev-smooth-approximation}
设 \(1\leq p<\infty\)、\(u\in W^{k,p}(\Omega)\)，且 \(u\) 在 \(\Omega\setminus K\) 上几乎处处为零，其中 \(K\Subset\Omega\)。对任意包含 \(K\) 的开集 \(V\subseteq\Omega\)，当 \(\varepsilon\) 充分小时，
\[
 u_\varepsilon=\rho_\varepsilon*\widetilde u\in C_c^\infty(V),
 \quad
 \|u_\varepsilon-u\|_{W^{k,p}(\Omega)}\longrightarrow0.
\]
更精确地，\(\operatorname{supp}u_\varepsilon\subseteq K+\overline B(0,\varepsilon)\)，其中右侧表示两集合中向量的所有和。
\end{proposition}
\begin{proof}
由内部紧支集补零引理，\(\widetilde u\in W^{k,p}(\mathbb R^d)\)，故前一命题给出全空间范数收敛，限制到 \(\Omega\) 后仍收敛。若 \(\operatorname{dist}(x,K)>\varepsilon\)，卷积所读取的点都在 \(K\) 外，因此卷积为零。由其连续性，支集包含于紧集 \(K+\overline B(0,\varepsilon)\)。当 \(K\ne\varnothing\) 时，取
\(\varepsilon<\operatorname{dist}(K,V^c)\)，便使这个紧集包含于 \(V\)；空补集的距离仍约定为无穷。若 \(K=\varnothing\)，原函数及其卷积均为零，结论直接成立。
\end{proof}

这个命题既控制逼近误差，也允许预先指定一个支集必须留在其中的开集。仅有误差趋零还不足以拼接局部逼近：支集若在磨光后聚集到同一点，原来局部有限的函数族就可能失去局部有限性。

\subsection{Meyers–Serrin 定理}
\label{b1:subsec:220103}

\term[meyers-serrin-dingli]{Meyers--Serrin 定理}[Meyers--Serrin theorem]把内部磨光推广到任意开集上的全域范数逼近。它要求逼近函数在 \(\Omega\) 内光滑，不要求这些函数可以光滑延伸到边界外，也不要求它们具有紧支集。

\begin{theorem}[Meyers--Serrin]\label{b1:thm:meyers-serrin}
设 \(\Omega\subseteq\mathbb R^d\) 为非空开集，\(k\in\mathbb N_0\)、\(1\leq p<\infty\)。则 \(C^\infty(\Omega)\cap W^{k,p}(\Omega)\) 在实或复空间 \(W^{k,p}(\Omega)\) 中稠密。

更一般地，给定 \(u\in W^{k,p}_{\mathrm{loc}}(\Omega)\) 及 \(\eta>0\)，存在 \(v\in C^\infty(\Omega)\)，使 \(v-u\in W^{k,p}(\Omega)\) 且
\[
 \|v-u\|_{W^{k,p}(\Omega)}<\eta.
\]
后一陈述不要求 \(u\) 或 \(v\) 各自属于全局 \(W^{k,p}(\Omega)\)。
\end{theorem}
\begin{proof}
证明较一般的陈述。沿用\cref{b1:thm:smooth-partition-unity}证明中的紧集穷竭，记
\[
 K_j=\{\,x\in\Omega\mid |x|\leq j,
                 \operatorname{dist}(x,\Omega^c)\geq1/j\,\},
 \quad K_0=K_{-1}=\varnothing,
\]
并置 \(V_j=\operatorname{int}K_{j+1}\setminus K_{j-2}\)。紧层 \(K_j\setminus\operatorname{int}K_{j-1}\) 覆盖 \(\Omega\)，且各自包含于 \(V_j\)，故这些开集也覆盖 \(\Omega\)。它们局部有限：任意一点有一个包含在某个 \(\operatorname{int}K_N\) 中的邻域，所有 \(j\geq N+2\) 的 \(V_j\) 都避开它。

该单位分解证明在每一层构造有限多个函数，其支集紧含于 \(V_j\)。把这些函数排成 \((\psi_i)_{i\geq1}\)，略去零函数，记第 \(i\) 个函数所属层为 \(j(i)\)。于是
\[
 \psi_i\geq0,
 \quad\sum_i\psi_i=1,
 \quad\operatorname{supp}\psi_i\Subset V_{j(i)},
\]
且每层只出现有限次。令 \(w_i=\psi_i u\)，并在 \(\Omega\) 外补零。前面的乘积讨论及补零引理保证 \(\widetilde w_i\in W^{k,p}(\mathbb R^d)\)。

对每个 \(i\)，使用\cref{b1:prop:compact-sobolev-smooth-approximation}，选取半径 \(\varepsilon_i>0\)，使
\[
 v_i=\rho_{\varepsilon_i}*\widetilde w_i\in C_c^\infty\bigl(V_{j(i)}\bigr),
 \quad
 \|v_i-w_i\|_{W^{k,p}(\Omega)}<\eta2^{-i-1}.
\]
这里先要求 \(\varepsilon_i<\operatorname{dist}\bigl(\operatorname{supp}\psi_i,V_{j(i)}^c\bigr)\)，再进一步缩小半径以达到误差条件。两个要求可以同时实现。

函数族 \((v_i)\) 仍局部有限。事实上，每个紧集只遇到有限层 \(V_j\)，而每层只有有限项。因此
\[
 v=\sum_{i=1}^\infty v_i
\]
在每个点的某个邻域内都是有限个光滑函数的和，故 \(v\in C^\infty(\Omega)\)。

剩下的全域估计应对误差求和。记 \(e_i=v_i-w_i\)。因为
\[
 \sum_i\|e_i\|_{W^{k,p}(\Omega)}\leq\eta/2,
\]
由\cref{b1:thm:sobolev-complete}及\ref{b1:prop:banach-series-criterion}，级数 \(\sum_i e_i\) 在 \(W^{k,p}(\Omega)\) 中收敛到某个 \(e\)，且 \(\|e\|_{W^{k,p}}\leq\eta/2\)。在任意内部相对紧开集上，\(v_i\) 和 \(w_i\) 都只有有限个非零项；足够长的部分和因而几乎处处等于
\[
 \sum_i(v_i-w_i)=v-u\sum_i\psi_i=v-u.
\]
另一方面，部分和在全域 \(L^p\) 中趋于 \(e\)，从而也在这个开集上趋于 \(e\)。故 \(e=v-u\) 几乎处处于该开集。再取可数个这样的开集穷竭 \(\Omega\)，得到全域等式与所需误差界。

若原来的 \(u\in W^{k,p}(\Omega)\)，则 \(v=u+e\) 也属于该空间。对每个正整数 \(n\) 取 \(\eta=1/n\)，便得到定理中的稠密性。
\end{proof}

证明没有宣称 \(\sum_{i=1}^N\psi_i u\) 在全域 Sobolev 范数中趋于 \(u\)。单位分解的导数可能随位置增大；局部有限只保证局部运算合法，并不自动给出全域范数收敛。逐项磨光后的误差另有可和估计，完备性用在这组误差上。

\subsection{\texorpdfstring{\(C_c^\infty\)}{Cc∞} 的稠密性问题}
\label{b1:subsec:220104}

Meyers--Serrin 定理产生的光滑函数可能一直延伸到边界附近或无穷远处。紧支集是额外要求，需要单独判断。

\begin{theorem}\label{b1:thm:whole-space-sobolev-density}
对 \(k\in\mathbb N_0\)、\(1\leq p<\infty\)，空间 \(C_c^\infty(\mathbb R^d)\) 在 \(W^{k,p}(\mathbb R^d)\) 中稠密。
\end{theorem}
\begin{proof}
由光滑截断引理，取 \(0\leq\chi\leq1\)，使 \(\chi=1\) 于 \(\overline B(0,1)\) 的一个邻域，且 \(\chi\in C_c^\infty\bigl(B(0,2)\bigr)\)。对 \(R\geq1\)，置 \(\chi_R(x)=\chi(x/R)\)。任取 \(u\in W^{k,p}(\mathbb R^d)\)。由 Leibniz 公式，对 \(|\alpha|\leq k\)，
\[
 \begin{aligned}
 \partial^\alpha(\chi_Ru-u)
 &=(\chi_R-1)\partial^\alpha u\\
 &\quad+\sum_{\substack{\beta\leq\alpha\\|\beta|>0}}
   \binom\alpha\beta R^{-|\beta|}
   (\partial^\beta\chi)(x/R)\partial^{\alpha-\beta}u.
 \end{aligned}
\]
第一项的 \(L^p\) 范数不超过
\(\|\mathbf1_{\mathbb R^d\setminus B(0,R)}\partial^\alpha u\|_p\)，由有限 \(p\) 的可积尾部趋于零。余下每项的范数不超过
\[
 \binom\alpha\beta R^{-|\beta|}
 \|\partial^\beta\chi\|_\infty\cdot\|\partial^{\alpha-\beta}u\|_p,
\]
也趋于零。有限个多重指标合起来给出 \(\chi_Ru\to u\) 于 \(W^{k,p}(\mathbb R^d)\)。

给定 \(\eta>0\)，先取 \(R\) 使截断误差小于 \(\eta/2\)。函数 \(\chi_Ru\) 具有紧支集，再由\cref{b1:prop:compact-sobolev-smooth-approximation}取 \(v\in C_c^\infty(\mathbb R^d)\)，使 \(\|v-\chi_Ru\|_{W^{k,p}}<\eta/2\)。三角不等式完成逼近。
\end{proof}

这里的截断只切掉无穷远处的尾部，并且 \(\chi_R\) 的非零阶导数带有衰减因子 \(R^{-|\beta|}\)。若试图在一般域的边界附近作越来越陡的截断，其导数反而可能增大，不能原样复制这个估计。

例如在 \(I=(0,1)\) 上取 \(u=1\)。它属于所有有限 \(p\) 的 \(W^{1,p}(I)\)。对任意 \(\varphi\in C_c^\infty(I)\)，函数 \(1-\varphi\) 的连续代表在两个端点都等于 \(1\)，因而由\eqref{b1:eq:interval-sobolev-sup}，
\[
 \begin{aligned}
 1&\leq\|1-\varphi\|_{L^1(I)}+\|\varphi'\|_{L^1(I)}\\
  &\leq\|1-\varphi\|_{L^p(I)}+\|\varphi'\|_{L^p(I)}\\
  &\leq2^{1-1/p}\|1-\varphi\|_{W^{1,p}(I)}.
 \end{aligned}
\]
最后一步是两项有限序列的 Hölder 不等式，也包括 \(p=1\)。因此这些紧支集函数到常数 \(1\) 的距离至少为 \(2^{1/p-1}>0\)。同一估计还排除了任意 \(k\geq1\) 时在 \(W^{k,p}(I)\) 中的逼近，因为高阶范数也控制其中的零阶和一阶项。问题来自边界附近强制为零的要求，不是域内光滑性不足；常数 \(1\) 本身已经光滑。

零阶情形应另作区分。对任意开集 \(\Omega\) 和有限 \(p\)，\(C_c^\infty(\Omega)\) 在 \(L^p(\Omega)\) 中仍稠密。事实上，取内部紧集穷竭 \(K_j\) 及 \(0\leq\chi_j\leq1\)、\(\chi_j\in C_c^\infty(\Omega)\)，使 \(\chi_j=1\) 于 \(K_j\)。则 \(\chi_j u\to u\) 于 \(L^p\)，这是控制收敛定理的直接应用；再对每个 \(\chi_j u\) 使用 \(k=0\) 的内部紧支集磨光命题即可。只控制函数值时，边界层的 \(L^p\) 误差可以消失；同时控制导数后，这个边界层便可能留下不可忽略的代价。下一节将把能够由紧支集光滑函数逼近的 Sobolev 函数单独组成一个闭子空间。
